%%%%%%%%%%%%%%%%%%%%%%%%%%%%%%%%%%%%%%%%%%%%%%%%%%%%%%%%%%%%%%
%%%%%%%%%%%%%%%%%% MA-0512 Teoría de conjuntos %%%%%%%%%%%%%%%%%%%%%%%%%%%%%%
%%%%%%%%%%%%%%%%%%%%%%%%%%%%%%%%%%%%%%%%%%%%%%%%%%%%%%%%%%%%%%
\newpage
\subsection*{MA-0512 Teoría de conjuntos}
\addcontentsline{toc}{subsection}{MA-0512 Teoría de conjuntos}

\hypertarget{MA0512}{}
\begin{refsection}
\begin{table}[H]
\centering{
\begin{tabular}{|ll|}
\hline
\textbf{Año:} IV  & \textbf{Requisitos:} Ninguno \\ 
\textbf{Ciclo:} Optativa  & \textbf{Correquisitos:} Ninguno \\ 
\textbf{Tipo de curso:} Teórico & \textbf{Horas presenciales por semana:} 5 \\ 
\textbf{Créditos:} 5 & \textbf{Horas de trabajo independiente por semana:} 10 \\
\textbf{Virtualidad:} P  &  \\ \hline
\end{tabular}
}
\end{table}

\subsubsection*{Descripción del curso}

Este curso optativo constituye una introducción a la teoría de conjuntos desde un enfoque axiomático. Se asume que el estudiantado posee nociones básicas sobre conjuntos y sus operaciones, así como experiencia con demostraciones matemáticas adquirida a lo largo de su formación. No se requiere un curso previo de lógica formal: el curso se desarrolla en el lenguaje matemático habitual, prestando atención al uso riguroso de cuantificadores y conectivos conforme estos aparecen de manera natural en los contenidos.

A lo largo del curso se desarrollará la teoría axiomática de Zermelo--Fraenkel, incluyendo el axioma de elección, con el propósito de proveer los fundamentos sobre los cuales se construye la matemática moderna. Se estudiarán los números cardinales y ordinales, el teorema de Cantor--Bernstein--Schroeder, la aritmética transfinita y resultados fundamentales sobre el continuo.

\subsubsection*{Objetivos}

\textbf{General:} Introducir al estudiantado en la teoría axiomática de conjuntos de Zermelo--Fraenkel, desarrollando las herramientas fundamentales como cardinales, ordinales y aritmética transfinita.

\textbf{Específicos:} Al finalizar el curso se espera que la persona estudiante sea capaz de:

\begin{enumerate}
\item Comprender y manejar los axiomas del sistema de Zermelo--Fraenkel con el axioma de elección (ZFC) como fundamento de la teoría de conjuntos.
\item Definir y analizar relaciones de equivalencia y de orden, incluyendo órdenes parciales, lineales y buenos órdenes, y aplicarlos en distintos contextos matemáticos.
\item Comprender el concepto de función como relación entre conjuntos y estudiar sus propiedades desde un punto de vista conjuntista.
\item Entender la noción de cardinalidad de un conjunto, distinguir entre conjuntos finitos, contables e incontables, y demostrar resultados fundamentales como el teorema de Cantor--Bernstein--Schroeder.
\item Desarrollar la teoría de números ordinales, incluyendo la aritmética ordinal y el principio de inducción transfinita, y aplicarlos en demostraciones matemáticas.
\item Estudiar la aritmética cardinal y su relación con el axioma de elección, la hipótesis del continuo y sus consecuencias.
\item Analizar formas equivalentes del axioma de elección, como el lema de Zorn y el teorema del buen orden, y discutir su rol en la matemática.
\item Demostrar de forma rigurosa los principales resultados del curso, comunicando las ideas matemáticas con claridad y precisión.
\end{enumerate}

\subsubsection*{Contenidos}

\begin{enumerate}

\item \textbf{Motivación y antecedentes históricos:}
\begin{enumerate}
    \item Breve historia de los fundamentos de la matemática.
    \item La paradoja de Russell y la necesidad de un enfoque axiomático.
    \item El programa de Hilbert y sus limitaciones.
    \item Panorama general de la teoría de conjuntos axiomática.
\end{enumerate}

\item \textbf{El sistema axiomático de Zermelo--Fraenkel:}
\begin{enumerate}
    \item El lenguaje de la teoría de conjuntos.
    \item Axioma de extensionalidad; igualdad de conjuntos.
    \item Axioma de separación (o especificación); subconjuntos.
    \item Axioma de la unión y axioma del par.
    \item Axioma del conjunto potencia.
    \item Axioma de reemplazo.
    \item Axioma de fundación (regularidad).
    \item Axioma del infinito.
\end{enumerate}

\item \textbf{Relaciones y funciones:}
\begin{enumerate}
    \item Pares ordenados y producto cartesiano.
    \item Relaciones: dominio, rango, imagen, composición e inversa.
    \item Relaciones de equivalencia y particiones.
    \item Funciones inyectivas, sobreyectivas y biyectivas.
    \item Operaciones con funciones; el conjunto de funciones entre dos conjuntos.
\end{enumerate}

\item \textbf{Órdenes:}
\begin{enumerate}
    \item Relaciones de orden parcial y de orden lineal.
    \item Elementos mínimos, minimales, máximos y maximales; cotas.
    \item Buen orden; el teorema del buen orden.
    \item Inducción sobre un buen orden.
    \item El lema de Zorn y sus aplicaciones.
\end{enumerate}

\item \textbf{Los números naturales:}
\begin{enumerate}
    \item El sucesor de un conjunto; conjuntos inductivos.
    \item El axioma del infinito y la existencia de $\omega$.
    \item Los axiomas de Peano como consecuencia de ZFC.
    \item Recursión sobre $\omega$; definición de las operaciones aritméticas.
    \item El orden en $\omega$; principio de inducción matemática.
\end{enumerate}

\item \textbf{Números cardinales:}
\begin{enumerate}
    \item Equipotencia de conjuntos; conjuntos finitos e infinitos.
    \item El teorema de Cantor--Bernstein--Schroeder.
    \item Conjuntos contables y no contables; el argumento diagonal de Cantor.
    \item El teorema de Cantor: $|A| < |\mathcal{P}(A)|$.
    \item Cardinales: definición y representantes canónicos.
    \item Aritmética cardinal: suma, producto y potencia de cardinales.
    \item Cardinales infinitos; el cardinal $\aleph_0$ y los cardinales $\aleph_\alpha$.
\end{enumerate}

\item \textbf{Números ordinales:}
\begin{enumerate}
    \item Conjuntos transitivos; definición de número ordinal.
    \item Propiedades básicas de los ordinales; el orden entre ordinales.
    \item Tipos de ordinales: sucesor y límite.
    \item Inducción transfinita y recursión transfinita.
    \item Aritmética ordinal: suma, producto y potencia.
    \item Relación entre ordinales y cardinales; ordinales iniciales.
\end{enumerate}

\item \textbf{El axioma de elección y la hipótesis del continuo:}
\begin{enumerate}
    \item El axioma de elección (AC): enunciado y formulaciones equivalentes.
    \item Equivalencia entre AC, el lema de Zorn y el teorema del buen orden.
    \item Consecuencias del axioma de elección en matemática.
    \item La hipótesis del continuo (CH) y la hipótesis generalizada del continuo (GCH).
    \item Los resultados de Gödel y Cohen: independencia de AC y CH respecto a ZF.
\end{enumerate}

\end{enumerate}

\subsubsection*{Metodología}

\noindent \textbf{Lineamientos metodológicos:} La persona docente a cargo de este curso debe respetar las siguientes pautas:

\begin{enumerate}
    \item Desarrollar el curso principalmente mediante \textbf{clases magistrales presenciales}, en las que se introduzcan los conceptos fundamentales, los resultados principales y las demostraciones más relevantes de la teoría, con participación activa del estudiantado.
    \item Complementar con \textbf{sesiones de ejercicios} para consolidar los conceptos teóricos, desarrollar habilidades de demostración y fortalecer el razonamiento lógico y abstracto.
    \item Abordar algunos temas mediante \textbf{talleres o proyectos} (individuales o grupales) orientados a explorar aplicaciones, profundizar en resultados específicos o estudiar desarrollos históricos y contemporáneos de la teoría de conjuntos.
    \item Promover evaluación formativa y sumativa; se sugiere incluir al menos dos exámenes parciales y un proyecto.
    \item Cuando las autoridades sanitarias o institucionales impongan restricciones, adaptar las actividades según normativa vigente.
\end{enumerate}

\textbf{Sugerencias metodológicas:} Adicionalmente, se tienen las siguientes recomendaciones para la persona docente a cargo de este curso:

\begin{enumerate}
    \item Utilizar \cite{Hal74} y \cite{End77} como referencias principales del curso, y \cite{Jec03} y \cite{Kun80} para temas avanzados como la independencia de AC y CH.
    \item Complementar con \cite{Dev93}, \cite{Sup72} y \cite{Hal17}; para aspectos de lógica, consultar \cite{HiL19}, \cite{CL2000} y \cite{EFT94}.
    \item Promover el aprendizaje autónomo, la investigación bibliográfica y la comunicación de ideas matemáticas de manera clara y rigurosa por medio de los proyectos.
\end{enumerate}

\nocite{Hal74}
\nocite{Jec03}
\nocite{Kun80}
\nocite{End77}
\nocite{Dev93}
\nocite{Sup72}
\nocite{Hal17}
\nocite{HiL19}
\nocite{CL2000}
\nocite{EFT94}

\printcursosbibliography
\end{refsection}
